% ============================================================
%  Příprava na maturitu – Mechatronika
%  Kompiluj: pdflatex maturita_mechatronika.tex
%  Kódování: UTF-8
% ============================================================

\documentclass[a4paper,10pt]{article}

% --- Balíčky ---
\usepackage[T1]{fontenc}
\usepackage[utf8]{inputenc}
\usepackage[czech]{babel}
\usepackage[left=2.4cm, right=1.8cm, top=1.8cm, bottom=1.8cm]{geometry}
\usepackage{lmodern}
\usepackage{microtype}
\usepackage{parskip}
\usepackage{titlesec}
\usepackage{enumitem}
\usepackage{xcolor}
\usepackage{hyperref}
\usepackage{tabularx}
\usepackage{booktabs}
\usepackage{array}
\usepackage{tikz}
\usepackage{eso-pic}
\usepackage{mdframed}
\usepackage{amsmath}
\usepackage{amssymb}
\usepackage{pgfplots}
\pgfplotsset{compat=1.18}
\usetikzlibrary{arrows.meta, positioning, shapes.geometric, calc}

% --- Barvy ---
\definecolor{accent}{HTML}{03254E}
\definecolor{stripeblue}{HTML}{568EA3}
\definecolor{medgray}{HTML}{888888}
\definecolor{lightblue}{HTML}{EAF2F8}
\definecolor{lightyellow}{HTML}{FDFBE4}
\definecolor{lightgreen}{HTML}{EAF7EA}

% --- Modrý pruh vlevo ---
\AddToShipoutPictureBG{%
  \begin{tikzpicture}[remember picture, overlay]
    \fill[stripeblue]
      (current page.north west)
      rectangle
      ([xshift=10mm] current page.south west);
    \fill[accent!60]
      ([xshift=10mm] current page.north west)
      rectangle
      ([xshift=12mm] current page.south west);
  \end{tikzpicture}%
}

\hypersetup{colorlinks=true, urlcolor=accent, linkcolor=accent}

\titleformat{\section}
  {\large\bfseries\color{accent}}
  {\thesection.}
  {0.5em}
  {}
  [\color{accent}\titlerule]

\titlespacing{\section}{0pt}{16pt}{6pt}

\newmdenv[
  backgroundcolor=lightblue,
  linecolor=stripeblue,
  linewidth=1pt,
  roundcorner=4pt,
  innerleftmargin=8pt,
  innerrightmargin=8pt,
  innertopmargin=6pt,
  innerbottommargin=6pt,
  skipabove=4pt,
  skipbelow=4pt
]{pojmybox}

\newmdenv[
  backgroundcolor=lightyellow,
  linecolor=accent!40,
  linewidth=1pt,
  roundcorner=4pt,
  innerleftmargin=8pt,
  innerrightmargin=8pt,
  innertopmargin=6pt,
  innerbottommargin=6pt,
  skipabove=4pt,
  skipbelow=8pt
]{vysvetlenibox}

\newmdenv[
  backgroundcolor=lightgreen,
  linecolor=accent!30,
  linewidth=1pt,
  roundcorner=4pt,
  innerleftmargin=8pt,
  innerrightmargin=8pt,
  innertopmargin=6pt,
  innerbottommargin=6pt,
  skipabove=4pt,
  skipbelow=8pt
]{vzorcebox}

\pagestyle{plain}

% ============================================================
\begin{document}

% --- Záhlaví ---
\begin{minipage}[t]{0.75\textwidth}
    {\Huge\bfseries\color{accent} Příprava na maturitu}\\[4pt]
    {\large\color{stripeblue} Mechatronika}\\[2pt]
    \textcolor{medgray}{\small Suchánek Lukáš \quad SPŠ Prosek \quad 2026}
\end{minipage}
\hfill
\begin{minipage}[t]{0.22\textwidth}
    \raggedleft
    \begin{tikzpicture}
        \draw[color=stripeblue, rounded corners=4pt, line width=1pt]
            (0,0) rectangle (3,1.2)
            node[midway, align=center, color=accent, font=\small\bfseries]
            {Otázek: 22};
    \end{tikzpicture}
\end{minipage}

\vspace{6pt}
\noindent\color{accent}\rule{\linewidth}{1.5pt}
\vspace{2pt}

% ============================================================
\section{Soustavy, statické, dynamické a frekvenční charakteristiky}

\begin{pojmybox}
\textbf{Klíčové pojmy:}
přenosová funkce, řádová soustava, přechodová charakteristika,
impulsní odezva, Bodeův diagram, frekvenční charakteristika,
fázová bezpečnost, amplitudová bezpečnost, póly, nuly
\end{pojmybox}

\begin{vysvetlenibox}
\textbf{Přenosová funkce}\\
Přenosová funkce je matematický model, který popisuje vstupně-výstupní
vztah lineární soustavy v komplexní oblasti. Používá se Laplaceova
transformace, která převádí diferenciální rovnice na algebraické rovnice.

\begin{vzorcebox}
$$G(s) = \frac{Y(s)}{U(s)} = \frac{b_m s^m + \cdots + b_0}
{a_n s^n + \cdots + a_0}$$
kde $s = j\omega$ je komplexní proměnná (Laplaceova transformace).\\
Soustava 1. řádu: $G(s) = \frac{K}{Ts + 1}$ \quad
(zesílení $K$, časová konstanta $T$)\\
Soustava 2. řádu: $G(s) = \frac{\omega_n^2}{s^2 + 2\zeta\omega_n s + \omega_n^2}$
\quad ($\zeta$ = poměrný útlum)
\end{vzorcebox}

\textbf{Póly a nuly přenosové funkce:}
\begin{itemize}[noitemsep, leftmargin=1.4em, topsep=4pt]
    \item \textbf{Nuly} -- kořeny čitatele; frekvence, při kterých je výstup nulový
    \item \textbf{Póly} -- kořeny jmenovatele; určují stabilitu a dynamiku
    \item \textbf{Stabilita:} všechny póly musí mít zápornou reálnou část
          (být v levé polorovině komplexní roviny)
\end{itemize}

\textbf{Typické soustavy podle řádu:}
\begin{itemize}[noitemsep, leftmargin=1.4em, topsep=4pt]
    \item \textbf{0. řád} -- $G(s) = K$; okamžitá odezva bez zpoždění
          (ideální zesilovač, dělič napětí)
    \item \textbf{1. řád} -- jeden pól; exponenciální náběh s časovou konstantou $T$
          (RC článek, ohřev s tělesem, nádrž s odtokem)
    \item \textbf{2. řád} -- dva póly; může kmitat při nízkém útlumu
          (pružina-hmota, RLC obvod, kyvadlo)
    \item \textbf{3. a vyšší řád} -- složené soustavy; často se aproximují
          na soustavu 1. nebo 2. řádu s dopravním zpožděním
\end{itemize}

\textbf{Dopravní zpoždění}\\
Některé systémy mají čisté zpoždění (materiál putuje pásem, signál se
zpracovává): $G(s) = G_0(s) \cdot e^{-Ls}$, kde $L$ je doba zpoždění.
Ve frekvenční oblasti: $e^{-j\omega L}$ -- fázový posun lineární s $\omega$.

\textbf{Charakteristiky soustav}

\begin{center}
\begin{tabularx}{0.95\linewidth}{@{} l X X @{}}
    \toprule
    \textbf{Charakteristika} & \textbf{Co popisuje} & \textbf{Jak se měří} \\
    \midrule
    Statická & Závislost výstupu na vstupu v ustáleném stavu;
               sklon = statické zesílení $K$
             & Nastavíme různé vstupy, čekáme na ustálení \\
    Přechodová & Odezva na jednotkový skok;
                 odečítáme $T$, $K$, přeregulaci $\sigma$
              & Skok vstupu, záznam výstupu \\
    Frekvenční & Bodeův diagram: $|G(j\omega)|$ a $\angle G(j\omega)$ vs. $\omega$
              & Harmonické buzení různých frekvencí \\
    Impulsní & Odezva na Diracův impuls; derivace přechodové
              & Krátký intenzivní zásah, záznam odezvy \\
    \bottomrule
\end{tabularx}
\end{center}

\textbf{Parametry přechodové charakteristiky:}
\begin{itemize}[noitemsep, leftmargin=1.4em, topsep=4pt]
    \item \textbf{$T_{63\%}$} -- čas, za který výstup dosáhne 63\% ustálené hodnoty
          (pro soustavu 1. řádu: $T_{63\%} = T$)
    \item \textbf{$T_{95\%}$} -- čas dosažení 95\% (pro 1. řád: $\approx 3T$)
    \item \textbf{Přeregulace $\sigma$} -- maximální překmit nad ustálenou hodnotu [\%]
    \item \textbf{Doba ustálení $T_u$} -- čas, kdy se výstup udrží v pásmu $\pm 5\,\%$
\end{itemize}

\textbf{Bodeův diagram -- princip}\\
Bodeův diagram se skládá ze dvou grafů:
\begin{itemize}[noitemsep, leftmargin=1.4em, topsep=4pt]
    \item \textbf{Amplitudová charakteristika} -- $|G(j\omega)|$ [dB] vs. $\omega$ [log]
    \item \textbf{Fázová charakteristika} -- $\angle G(j\omega)$ [°] vs. $\omega$ [log]
\end{itemize}

\textbf{Asymptotický průběh pro soustavu 1. řádu:}
\begin{itemize}[noitemsep, leftmargin=1.4em, topsep=4pt]
    \item Pro $\omega < \omega_c$ (lomová frekvence): vodorovná čára na 0\,dB
    \item Pro $\omega > \omega_c$: klesající čára se sklonem $-20$\,dB/dekádu
    \item Fáze: přechod od 0° k $-90°$ v okolí $\omega_c$
\end{itemize}

\begin{center}
\begin{tikzpicture}
\begin{semilogxaxis}[
  width=9cm, height=4cm,
  xlabel={Frekvence $\omega$ [rad/s]},
  ylabel={Zesílení [dB]},
  xmin=0.01, xmax=100,
  ymin=-40, ymax=5,
  grid=major, grid style={dotted,medgray},
  tick label style={font=\small},
  label style={font=\small},
  title={Bodeův diagram $G(s)=\frac{1}{s+1}$},
  title style={font=\small\bfseries}
]
\addplot[accent, thick, domain=0.01:100, samples=200]
    {-10*log10(1 + x^2)};
\addplot[stripeblue, dashed, domain=0.01:1] {0};
\addplot[stripeblue, dashed, domain=1:100]  {-20*(log10(x))};
\draw[red!70, dashed] (axis cs:1,-40) -- (axis cs:1,5)
    node[above, font=\scriptsize, color=red!70] {$\omega_c=1/T$};
\end{semilogxaxis}
\end{tikzpicture}
\end{center}

\textbf{Pravidla pro kreslení Bodeova diagramu:}
\begin{itemize}[noitemsep, leftmargin=1.4em, topsep=4pt]
    \item \textbf{Zesílení $K$} -- posune amplitudu o $20\log_{10}K$ dB
    \item \textbf{Pól v počátku} ($1/s$) -- sklon $-20$\,dB/dek, fáze $-90°$
    \item \textbf{Pól reálný} ($1/(Ts+1)$) -- lom při $\omega=1/T$, sklon $-20$\,dB/dek
    \item \textbf{Dvojice pólů} (2. řád) -- sklon $-40$\,dB/dek, fáze $-180°$
    \item \textbf{Nula} -- opačný efekt než pól (kladný sklon)
\end{itemize}

\textbf{Stabilitní kritéria z Bodeova diagramu}

\begin{vzorcebox}
\textbf{Fázová bezpečnost:} $\varphi_m = 180° + \angle G(j\omega_0)$
kde $\omega_0$ je frekvence, při níž $|G|=0\,\text{dB}$.
Požadujeme $\varphi_m > 30°$ (typicky $45°$--$60°$).\\[4pt]
\textbf{Amplitudová bezpečnost:} $A_m = -|G(j\omega_{\pi})|$
kde $\omega_{\pi}$ je frekvence, při níž $\angle G = -180°$.
Požadujeme $A_m > 6\,\text{dB}$.
\end{vzorcebox}

\textbf{Význam bezpečností:}
\begin{itemize}[noitemsep, leftmargin=1.4em, topsep=4pt]
    \item \textbf{Fázová bezpečnost} -- jak daleko jsme od vzniku kmitů;
          malá $\varphi_m$ $\to$ velké překmity, oscilace
    \item \textbf{Amplitudová bezpečnost} -- jak můžeme zvýšit zesílení,
          než systém ztratí stabilitu
    \item \textbf{Typické hodnoty:} $\varphi_m = 45°\text{--}60°$, $A_m = 6\text{--}12$\,dB
\end{itemize}

\textbf{Nyquistův diagram}\\
Alternativa k Bodeovu diagramu -- komplexní přenosová funkce
$G(j\omega)$ se vykreslí do jedné křivky v komplexní rovině.
\begin{itemize}[noitemsep, leftmargin=1.4em, topsep=4pt]
    \item Vzdálenost od bodu $[-1, j0]$ určuje stabilitu
    \item Kritério: křivka nesmí obklopit bod $[-1, j0]$
    \item Výhoda: vidíme amplitudu i fázi najednou
\end{itemize}

\textbf{Identifikace soustavy z měření}\\
Postup, když neznáme matematický model:
\begin{enumerate}[noitemsep, leftmargin=1.4em, topsep=4pt]
    \item Změříme přechodovou charakteristiku (skok vstupu)
    \item Odečteme $T_{63\%}$, $T_{28\%}$, přeregulaci $\sigma$
    \item Typ soustavy: bez přeregulace $\to$ 1. řád; s přeregulací $\to$ 2. řád
    \item Vypočítáme parametry: $T = T_{63\%}$, $K = y(\infty)/u_{skok}$
    \item Ověříme na frekvenční charakteristice
\end{enumerate}
\end{vysvetlenibox}

% ============================================================

% ============================================================
\end{document}
